
Comme $i'$ est régulière, et $p$ et $p'$ lisses, on en déduit d'après 1.3 que $i''$ et $i$ sont régulières.

La proposition qui suit ne sera pas utilisée dans la suite du Séminaire :

\textbf{Proposition 1.4.} \uline{Soit} $f:X\to Y$ \uline{un morphisme plat localement de présentation finie}. \uline{Pour que} $f$ \uline{soit un morphisme d'intersection complète au sens de} 1.1, \uline{il faut et il suffit qu'il soit un morphisme d'intersection complète au sens de} EGA IV 19.3.6.

Supposons que $f$ soit un morphisme d'intersection complète au sens de 1.1, et soit $x\in X$. Alors il existe un voisinage ouvert $U$ de $x$, et une factorisation :
\[
\begin{tikzcd}[column sep=large,row sep=large]
U \arrow[rr,"i"] \arrow[dr,"f"'] && V \arrow[dl,"g"]\\
& Y &
\end{tikzcd}
\]
où $i$ est une immersion régulière, et $g$ un morphisme lisse. Soit $y=f(x)$ ; alors l'immersion $i_y:U_y\to V_y$ est régulière en $x$. En effet, on peut supposer, quitte à restreindre $V$, qu'il existe un $\underline{O}_V$-Module libre de type fini $\underline{E}$, et un homomorphisme régulier $u:\underline{E}\to\underline{O}_V$ tel que le complexe de Koszul associé à $u$ soit une résolution de $\underline{O}_U$. On considère alors l'homomorphisme $u_y$ induit sur la fibre en $y$, et le complexe de Koszul associé $\underline{K}_\bullet(u_y)=\underline{K}_\bullet(u)_y$. Comme $g$ est lisse, donc plat, $\underline{K}_\bullet(u)$ est une résolution $Y$-plate de $\underline{O}_U$. Par suite, les groupes d'homologie de $\underline{K}_\bullet(u)_y$ sont les $\underline{\Tor}_i^{\underline{O}_Y}(\underline{O}_U,\underline{k}(y))$, où $\underline{k}(y)$ est le corps résiduel de $\underline{O}_{Y,y}$. Comme $f$ est un morphisme plat, ils sont nuls, et $i_y$ est régulière en $x$. Par ailleurs, comme $g$ est lisse, l'anneau local de $V_y$ en $x$ est régulier. Comme l'anneau local de $U_y$ en $x$ est quotient d'un anneau régulier par un idéal régulier,
